\documentclass[11pt]{article}

\usepackage[a4paper,margin=1in]{geometry}
\usepackage{amsmath,amssymb,amsthm,mathtools}
\usepackage{bm}
\usepackage{booktabs}
\usepackage{enumitem}
\usepackage{cite}
\usepackage{hyperref}

\title{\bfseries Handles Before Interventions:\\
Access-Model UAD and the Embedded Semantics of Agency Tests}
\author{Gunnar Zarncke\\AE Studio}
\date{\today}

\newtheorem{definition}{Definition}
\newtheorem{assumption}{Assumption}
\newtheorem{proposition}{Proposition}
\newtheorem{theorem}{Theorem}
\newtheorem{corollary}{Corollary}
\newtheorem{remark}{Remark}

\newcommand{\X}{\mathcal X}
\newcommand{\Y}{\mathcal Y}
\newcommand{\U}{\mathcal U}
\newcommand{\Hh}{\mathcal H}
\newcommand{\A}{\mathcal A}
\newcommand{\D}{\mathcal D}
\newcommand{\M}{\mathcal M}
\newcommand{\C}{\mathcal C}
\newcommand{\E}{\mathcal E}
\newcommand{\I}{\mathcal I}
\newcommand{\G}{\mathcal G}
\newcommand{\Z}{\mathcal Z}
\newcommand{\Prob}{\mathbb P}
\newcommand{\Expect}{\mathbb E}
\newcommand{\KL}{D_{\mathrm{KL}}}
\newcommand{\MI}{I}
\newcommand{\doop}{\operatorname{do}}
\newcommand{\DL}{\operatorname{DL}}
\newcommand{\EIG}{\operatorname{EIG}}
\newcommand{\sem}{\operatorname{sem}}
\newcommand{\cost}{\operatorname{Cost}}
\newcommand{\risk}{\operatorname{Risk}}
\newcommand{\argmax}{\operatorname*{arg\,max}}
\newcommand{\argmin}{\operatorname*{arg\,min}}

\begin{document}
\maketitle

\begin{abstract}
Unsupervised Agent Discovery (UAD)~\cite{uad2025} is usually stated as if the observer receives
variables, partitions them into internal, sensory, active, and external components,
and then applies conditional-independence or control-information tests. This
paper argues that this picture is too shallow for real systems. In laboratories,
software audits, social systems, biological assays, recommender platforms, and
human interaction, neither ``intervention'' nor ``interface'' is primitive. What
the observer has are \emph{handles}: embedded access points through which the
system may be observed or operated upon. A handle has observational semantics,
operational semantics, cost, latency, reversibility, spillover, and uncertain
uptake. Causal graph interventions such as $\doop(X_i=x)$~\cite{pearl2009causality} are idealizations of
rare depth-zero cases. We introduce an access-model formulation of UAD in which
the data are generated through observation kernels and operation kernels. Agent
loops, interfaces, and interventions are then inferred jointly. The resulting
Bayesian/regularized objective draws on blanket-information competence scoring~\cite{biq2025}
and treats agency as posterior mass over handle-mediated
control-loop hypotheses, not as a label on a raw time series. We prove an
impossibility result: no method can distinguish agent structures that are
equivalent under the available access model. We then give a margin-style recovery
bound for finite candidate classes, an intervention-value criterion based on
expected information gain over agent-loop posteriors, and a sparse layer-wise
regularization scheme that encourages interpretable causal roles at every scale
while allowing full agency only at selected coarse-grained bottlenecks. The
framework explains why algorithmic feeds and advertising systems can manipulate
humans while hiding agency: users see content handles but lack access to objective,
ranking, telemetry, and policy-update handles. Interface forcing is therefore
not merely convenience; it is access-channel design.
\end{abstract}

\section{Motivation}

UAD begins from an appealing mathematical fantasy. There is a dynamical system
\[
    X_t=(X_1(t),\ldots,X_n(t)),
\]
and one seeks a subset $C\subseteq\{1,\ldots,n\}$ with an internal--sensory--active--external
decomposition
\[
    C = I^C\cup S^C\cup A^C,\qquad E^C = \{1,\ldots,n\}\setminus C,
\]
such that
\begin{equation}
    J(C) :=
    \MI\!\left(I^C_{t+1};E^C_{t+1}\mid S^C_t,A^C_t\right)
    \approx 0 .
    \label{eq:uad-blanket}
\end{equation}
The test says: conditional on sensory and active boundary variables, internal
and external futures are approximately screened off.

This is a useful formal target. But it hides a prior ontological assumption:
the observer has already been given the right variables at the right scale and
knows enough of the interface to call some variables sensory and some active.
Likewise, when we say ``intervene on $X_i$'', we often implicitly assume the
Pearlean ideal~\cite{pearl2009causality}
\[
    \doop(X_i=x),
\]
as if the observer could surgically set one coordinate of the system. Real
systems rarely offer such access. A chemical perturbation changes many pathways.
A prompt changes a distribution over latent computations. A platform audit flag
changes incentives, ranking policies, and future strategic behavior. A social
request is interpreted by another agent with its own beliefs. Even a laboratory
apparatus has uptake, delay, measurement error, spillover, and constraints.

Thus the primitive of intervention-capable UAD should not be a variable. It
should be a \emph{handle}. A handle is an access point with observational and/or
operational semantics. The observer does not perform $\doop(X_i=x)$; the
observer applies an operation $u$ through a handle $h$, producing a stochastic
change in future observations:
\[
    u@h \leadsto P(Y_{t+1:t+H}\mid Y_{\le t},u,h).
\]
The question becomes: given an access model, what agent loops are identifiable,
and which handle-operations should be used to test them?

\section{Worlds, Observations, and Handles}

Let $(X_t)_{t\ge0}$ be an unobserved world process with state space $\X$. The
observer does not see $X_t$ directly. Instead, observations arrive through a set
of handles.

\begin{definition}[Observation handle]
An observation handle is a tuple
\[
    h^Y = (\Y_h,K_h,\ell_h,\sigma_h,c_h),
\]
where $\Y_h$ is the handle's observation space,
\[
    K_h:\X\times \X^{-}\to \Delta(\Y_h)
\]
is an observation kernel possibly depending on current and recent world states,
$\ell_h$ is latency, $\sigma_h$ is an uncertainty or resolution descriptor, and
$c_h$ is the cost of using the handle. The observed datum is
\[
    Y^h_t \sim K_h(\cdot\mid X_{t-\ell_h:t}).
\]
\end{definition}

\begin{definition}[Operation handle]
An operation handle is a tuple
\[
    h^U = (\U_h,\Gamma_h,\rho_h,\chi_h,c_h),
\]
where $\U_h$ is a space of available operations,
\[
    \Gamma_h:\X\times \U_h\to \Delta(\X)
\]
is an operation kernel, $\rho_h$ is reversibility, $\chi_h$ describes constraints
or legality, and $c_h:\U_h\to\mathbb R_{\ge0}$ is operation cost. Applying
operation $u\in\U_h$ at time $t$ induces
\[
    X_{t+1}\sim \Gamma_h(\cdot\mid X_t,u)
\]
or, more generally, a finite-horizon kernel
\[
    X_{t+1:t+H}\sim \Gamma_h^H(\cdot\mid X_{\le t},u).
\]
\end{definition}

\begin{definition}[Access model]
An access model is
\[
    \A
    =
    (\Hh^Y,\Hh^U,K,\Gamma,\cost,\risk,\mathcal C),
\]
where $\Hh^Y$ is a family of observation handles, $\Hh^U$ a family of operation
handles, $K=\{K_h\}_{h\in\Hh^Y}$, $\Gamma=\{\Gamma_h\}_{h\in\Hh^U}$, $\cost$
and $\risk$ assign costs and risks to handle-use, and $\mathcal C$ encodes
constraints on allowable observations and operations.
\end{definition}

The data available to UAD are therefore not merely $X_{1:T}$ but
\[
    \D_T =
    \{(h_t^Y,y_t), (h_t^U,u_t), \text{context}_t\}_{t=1}^T,
\]
together with partial knowledge or priors over $K$ and $\Gamma$.

\begin{remark}
The passive time-series case is recovered by taking $\Hh^U=\varnothing$ and a
single observation handle $h^Y$ with $Y_t=X_t+\eta_t$. The classical UAD setting~\cite{uad2025,ConantAshby1970,kirchhoff2018markov}
is therefore a special, high-access, low-embedding case.
\end{remark}

\section{From Variables to Handles}

In classical causal notation~\cite{pearl2009causality}, intervention is represented as graph surgery:
\[
    P(X_{-i}\mid \doop(X_i=x)).
\]
The handle view replaces this with an operation kernel:
\[
    P(Y_{t+1:t+H}\mid \D_{\le t},u,h).
\]
The relationship between the two is approximate and itself inferential.

\begin{definition}[Effective target distribution]
For operation $u$ through handle $h$, define the effective target posterior
\[
    \tau_{h,u}(i)
    :=
    \Prob\!\left(
        X_i \text{ is directly affected by }(h,u)
        \mid \D
    \right).
\]
The spillover posterior is
\[
    \varsigma_{h,u}(j)
    :=
    \Prob\!\left(
        X_j \text{ is indirectly affected by }(h,u)
        \mid \D
    \right).
\]
\end{definition}

An ideal intervention is one with concentrated $\tau$ and low $\varsigma$:
\[
    H(\tau_{h,u})\approx 0,
    \qquad
    \|\varsigma_{h,u}\|_1-\max_i\tau_{h,u}(i)\approx0.
\]
Most real interventions have broad target and spillover distributions.

\begin{definition}[Intervention depth]
The depth of an operation $u@h$ relative to a hypothesized target variable $X_i$
is the minimal causal path length, in the observer's current structural posterior,
from the operation event $U_t=(h,u)$ to $X_i$.
Depth zero corresponds to $\doop(X_i=x)$. Depth one corresponds to a known
actuator on $X_i$. Depth two or greater corresponds to mediated perturbations
through apparatus, environment, norms, incentives, language, or other agents.
\end{definition}

The depth notion is not intended as a fixed graph invariant. It is posterior and
access-relative:
\[
    P(\operatorname{depth}(u@h,X_i)=d\mid \D).
\]

\section{Interfaces as Coupling Surfaces}

A sensor/action interface is not generally a set of raw variables. It is a stable
coupling surface, often constructed by an observer.

Let $C$ be a candidate system and $E$ its complement at some representation
level. A candidate interface $Q$ is a handle-mediated representation of coupling
between $C$ and $E$.

\begin{definition}[Candidate interface]
A candidate interface for $C$ is
\[
    Q=(Q^S,Q^A,K_Q,\Gamma_Q),
\]
where $Q^S$ are candidate sensory handles/channels, $Q^A$ candidate active
handles/channels, $K_Q$ describes how interface observations are produced, and
$\Gamma_Q$ describes how interface operations or outputs affect the world.
\end{definition}

The interface should approximately screen off internal and external futures:
\begin{equation}
    J_Q(C)
    :=
    \MI(C_{t+1};E_{t+1}\mid Q_t)
    \quad\text{small}.
    \label{eq:interface-screening}
\end{equation}
For directed sensor/action roles, define the asymmetric information scores
\begin{align}
    R_S(Q;C)
    &=
    \MI(E_t;Q^S_t\mid C_t)
    +
    \MI(Q^S_t;C_{t+1}\mid C_t),
    \\
    R_A(Q;C)
    &=
    \MI(C_t;Q^A_t\mid E_t)
    +
    \MI(Q^A_t;E_{t+1}\mid E_t).
\end{align}
These are not sufficient for agency, but they localize candidate inbound and
outbound channels.

\begin{definition}[Handle interface score]
For candidate $C$ and interface $Q$, define
\begin{equation}
    \mathcal I(C,Q)
    =
    -J_Q(C)
    +\alpha_S R_S(Q;C)
    +\alpha_A R_A(Q;C)
    -\lambda_Q \DL(Q)
    +\mu_Q \operatorname{Manip}(Q).
    \label{eq:interface-score}
\end{equation}
Here $\operatorname{Manip}(Q)$ measures whether available operation handles can
selectively affect or read the interface:
\[
    \operatorname{Manip}(Q)
    =
    \max_{h,u}
    \MI(u;Q_{t+1:t+H}\mid \D_{\le t},h).
\]
\end{definition}

\begin{remark}
Interface forcing is the deliberate construction of a high-$\mathcal I(C,Q)$
access channel. An API, laboratory assay, telemetry schema, randomized audit
interface, or standardized social probe is not merely a measurement. It changes
the access model and may change the system.
\end{remark}

\section{Agent-Loop Hypotheses}

An agent-loop hypothesis is a structured explanation of a candidate bounded
process as a stateful control system.

\begin{definition}[Agent-loop hypothesis]
An agent-loop hypothesis is
\[
    \Omega
    =
    (C,Q,B,G,\pi,V,\Theta),
\]
where $C$ is a candidate boundary, $Q=(Q^S,Q^A)$ a candidate interface, $B_t$ a
latent belief or information state, $G_t$ a latent goal/preference/viability
state, $\pi$ a policy, $V$ a viability or objective functional, and $\Theta$
additional parameters. The model class induced by $\Omega$ predicts observations
and responses by
\begin{align}
    B_{t+1}
    &\sim
    U_\Theta(B_t,Q^S_{t+1},Q^A_t),
    \\
    G_{t+1}
    &\sim
    R_\Theta(G_t,B_t,Q_t),
    \\
    Q^A_t
    &\sim
    \pi_\Theta(\cdot\mid B_t,G_t),
    \\
    Y_{t+1}
    &\sim
    P_\Theta(\cdot\mid B_{t+1},G_{t+1},Q^A_t,E_t).
\end{align}
\end{definition}

The latent variables $B_t$ and $G_t$ need not be internal physical states. They
are explanatory coordinates under which behavior compresses and interventions
become predictable.

\section{Bayesian Access-Model UAD}

The full posterior should include both agent-loop hypotheses and access semantics:
\[
    P(\Omega,K,\Gamma,Q\mid \D_T).
\]
This is the core object. It replaces both ``detect an agent in a time series''
and ``intervene on a variable''.

\begin{definition}[Access-model posterior]
Given priors over access semantics and loop hypotheses,
\begin{equation}
    P(\Omega,K,\Gamma,Q\mid \D_T)
    \propto
    P(\D_T\mid \Omega,K,\Gamma,Q)
    P(\Omega\mid Q)
    P(Q\mid K,\Gamma)
    P(K,\Gamma).
    \label{eq:posterior}
\end{equation}
\end{definition}

To make this computationally tractable, use an energy parameterization:
\[
    P(\Omega,Q\mid \D_T)
    \propto
    \exp[-E(\Omega,Q;\D_T)].
\]
A generic energy is
\begin{align}
    E(\Omega,Q;\D_T)
    &=
    \lambda_J J_Q(C)
    -\lambda_C I_{\mathrm{ctrl}}(\Omega)
    -\lambda_P R_{\mathrm{persist}}(\Omega)
    \nonumber\\
    &\qquad
    -\lambda_B R_{\mathrm{belief}}(\Omega)
    -\lambda_V R_{\mathrm{viability}}(\Omega)
    -\lambda_L \Delta L_{\mathrm{EIS}}(\Omega)
    \nonumber\\
    &\qquad
    -\lambda_M \operatorname{Manip}(Q)
    +\kappa_\Omega \DL(\Omega)
    +\kappa_Q \DL(Q)
    \nonumber\\
    &\qquad
    +\kappa_A \DL(K,\Gamma).
    \label{eq:energy}
\end{align}
where $\Delta L_{\mathrm{EIS}}$ is the intentional compression gain~\cite{eis2025}.

The terms have the following intended meanings:
\begin{align}
    I_{\mathrm{ctrl}}(\Omega)
    &=
    \MI(Q^A_t;E_{t+1:t+H}\mid E_t,Q^S_t),
    \\
    R_{\mathrm{persist}}(\Omega)
    &=
    \MI(G_t;G_{t+k})-\eta_G H(G_t),
    \\
    R_{\mathrm{belief}}(\Omega)
    &=
    \MI(Q^A_t;B_t\mid E_t,G_t),
    \\
    R_{\mathrm{viability}}(\Omega)
    &=
    \Expect[\Delta V_{t:t+H}\mid \pi_\Omega]
    -
    \Expect[\Delta V_{t:t+H}\mid \pi_{\mathrm{null}}],
    \\
    \Delta L_{\mathrm{EIS}}(\Omega)
    &=
    \log P(\D_T\mid M_{\mathrm{belief/goal}},\Omega)
    -
    \log P(\D_T\mid M_{\mathrm{null}},Q)
    -
    \lambda_{\mathrm{int}}\DL(M_{\mathrm{belief/goal}}).
\end{align}

\begin{remark}
The non-agentic baseline should not be an unconstrained universal predictor.
A powerful predictor can absorb agentic regularities without representing agent
loops. Instead, the null family should consist of constrained causal topologies:
passive dynamics, reactive dynamics, scripted dynamics, field dynamics, and
habitual dynamics without intervention-sensitive goals.
\end{remark}

\section{Null Models as Restricted Topologies}

Let
\[
    \M = \{M_A,M_0^1,\ldots,M_0^k\},
\]
where $M_A$ is an agent-loop model and $M_0^i$ are restricted alternatives. The
posterior agent probability is
\begin{equation}
    P(M_A\mid \D)
    =
    \frac{
        P(\D\mid M_A)P(M_A)
    }{
        P(\D\mid M_A)P(M_A)
        +
        \sum_i P(\D\mid M_0^i)P(M_0^i)
    }.
    \label{eq:model-posterior}
\end{equation}

Useful nulls include:
\begin{align}
    M_0^{\mathrm{passive}}&:\quad
    X_{t+1}=F(X_t)+\eta_t,
    \\
    M_0^{\mathrm{reactive}}&:\quad
    Q^A_t=f(Q^S_t)+\eta_t,
    \\
    M_0^{\mathrm{script}}&:\quad
    Q^A_t=f(t,\text{context})+\eta_t,
    \\
    M_0^{\mathrm{field}}&:\quad
    X_{t+1}=F_{\mathrm{global}}(X_t)+\eta_t,
    \\
    M_0^{\mathrm{habit}}&:\quad
    Q^A_t=f(Q^S_{\le t},Q^A_{<t})+\eta_t,
\end{align}
where none includes a persistent goal variable that remains stable under changed
means and predicts counterfactual policy adaptation.

\section{Layered Role Regularization}

It is too strong to require agentic structure at every representation layer.
Low-level features may be pixels, tokens, molecules, or local dynamical modes.
They need not be agentic. But they should become locally understandable.

Let a representation learner produce layers
\[
    Z^0_t=X_t,\qquad
    Z^\ell_t=f^\ell_\theta(Z^{\ell-1}_{\le t}).
\]
At each layer $\ell$, introduce soft role masks
\[
    m_C^\ell,m_Q^\ell,m_E^\ell,m_B^\ell,m_G^\ell \in [0,1]^{d_\ell}.
\]
Use role regularization at all layers:
\begin{align}
    R^\ell_{\mathrm{role}}
    &=
    \rho_1 R^\ell_{\mathrm{persistence}}
    +\rho_2 R^\ell_{\mathrm{locality}}
    +\rho_3 R^\ell_{\mathrm{separability}}
    +\rho_4 R^\ell_{\mathrm{causal-handle}}
    -\rho_5 R^\ell_{\mathrm{degenerate}}.
\end{align}
But introduce sparse agent-loop gates
\[
    a^\ell_C\in\{0,1\}
\]
and apply the full agent score only when $a^\ell_C=1$:
\begin{equation}
    \mathcal L
    =
    \mathcal L_{\mathrm{pred}}
    -
    \sum_\ell \lambda_{\mathrm{role}} R^\ell_{\mathrm{role}}
    -
    \sum_{\ell,C} q(a^\ell_C=1)\lambda_{\mathrm{agent}}R^\ell_{\mathrm{agent}}(C)
    +
    \tau\sum_{\ell,C}
    \KL\!\left(q(a^\ell_C)\,\|\,p(a^\ell_C)\right).
    \label{eq:layer-objective}
\end{equation}
The prior $p(a^\ell_C=1)$ is sparse. Most regions at most layers are not agents.

\begin{definition}[Agentic structure prior]
For a layer $\ell$ and candidate $C$,
\[
    p(\Omega^\ell_C)
    \propto
    \exp\left(
        R^\ell_{\mathrm{agent}}(C)
        -
        \kappa \DL(\Omega^\ell_C)
    \right).
\]
Thus agentic regularization is the log prior over latent loop structures.
\end{definition}

\section{Coarse-Graining and Hierarchical Agency}

Representation layers are not necessarily ontological layers. Hierarchy should
therefore be searched over coarse-graining operators
\[
    K_{\sigma,\tau}:\X_{1:T}\to \Z^{\sigma,\tau}_{1:T},
\]
where $\sigma$ denotes variable/spatial aggregation and $\tau$ temporal
aggregation. Define
\begin{equation}
    \operatorname{ScaleAgent}(C)
    =
    \Expect_{\sigma,\tau}
    \left[
        R_{\mathrm{agent}}(C;K_{\sigma,\tau})
    \right]
    -
    \lambda_{\mathrm{var}}
    \operatorname{Var}_{\sigma,\tau}
    \left[
        R_{\mathrm{agent}}(C;K_{\sigma,\tau})
    \right].
    \label{eq:scale-agent}
\end{equation}
A hierarchical agent is a stable basin in scale-space, not merely an activation
in a neural-network layer.

\section{Identifiability Under Access Models}

The access model determines what can be known.

\begin{definition}[Access equivalence]
Two agent-loop hypotheses $\Omega$ and $\Omega'$ are equivalent under access
model $\A$, written $\Omega\sim_\A\Omega'$, if for every admissible observation
and operation policy $\nu$ over handles,
\[
    P_\Omega(\D_T^\nu\mid \A)
    =
    P_{\Omega'}(\D_T^\nu\mid \A)
    \qquad
    \forall T.
\]
The recoverable target is the equivalence class
\[
    [\Omega]_\A=\{\Omega':\Omega'\sim_\A\Omega\}.
\]
\end{definition}

\begin{theorem}[Access-model impossibility]
If $\Omega\sim_\A\Omega'$, then no UAD procedure whose inputs are restricted to
data generated through $\A$ can distinguish $\Omega$ from $\Omega'$ with
probability better than chance under all priors assigning equal mass to
$\Omega$ and $\Omega'$.
\end{theorem}

\begin{proof}
By access equivalence, every admissible data-generating distribution is identical
under $\Omega$ and $\Omega'$. Hence the likelihood ratio is identically one:
\[
    \frac{P(\D_T\mid \Omega,\A)}{P(\D_T\mid \Omega',\A)}=1.
\]
With equal prior mass, posterior odds remain one for all $T$ and all admissible
operation policies. Any decision rule must therefore have symmetric error.
\end{proof}

\begin{remark}
This theorem is the handle-level analogue of observation-channel
non-identifiability. If the access model hides the objective, policy update,
telemetry, or actuator channels of a platform, then user-level observations may
not distinguish an optimizing recommender from a non-agentic content stream
with matching distribution.
\end{remark}

\section{Finite Candidate Recovery}

Let $\mathcal O$ be a finite class of candidate loop/interface hypotheses.
Let $S_\A(\Omega)$ be a population score, for example the negative energy in
Eq.~\eqref{eq:energy}, and $\widehat S_\A(\Omega)$ its empirical estimate.

\begin{assumption}[True recoverable class]
There is a true access-equivalence class $[\Omega^\star]_\A$.
\end{assumption}

\begin{assumption}[Access margin]
For all $\Omega\notin[\Omega^\star]_\A$,
\[
    S_\A(\Omega^\star)-S_\A(\Omega)\ge \Delta_\A>0.
\]
\end{assumption}

\begin{assumption}[Access distortion]
The implemented handles and estimated kernels distort the ideal access score by
at most
\[
    \sup_{\Omega\in\mathcal O}
    |S_{\widehat\A}(\Omega)-S_\A(\Omega)|
    \le \delta_\A.
\]
\end{assumption}

\begin{assumption}[Uniform concentration]
For all $\epsilon>0$,
\[
    \Prob\left[
        \sup_{\Omega\in\mathcal O}
        |\widehat S_{\widehat\A}(\Omega)-S_{\widehat\A}(\Omega)|
        >
        \epsilon
    \right]
    \le
    2|\mathcal O|
    \exp\left(
        -\frac{T_{\mathrm{eff}}\epsilon^2}{2B^2}
    \right).
\]
\end{assumption}

\begin{theorem}[Handle-recovery bound]
Let
\[
    \widehat\Omega
    \in
    \argmax_{\Omega\in\mathcal O}
    \widehat S_{\widehat\A}(\Omega).
\]
If
\[
    \Delta_\A>2\delta_\A,
\]
then
\[
    \Prob\left[
        \widehat\Omega\in[\Omega^\star]_\A
    \right]
    \ge
    1
    -
    2|\mathcal O|
    \exp\left(
        -\frac{T_{\mathrm{eff}}(\Delta_\A-2\delta_\A)^2}{8B^2}
    \right).
\]
\end{theorem}

\begin{proof}
For any $\Omega\notin[\Omega^\star]_\A$,
\[
    S_{\widehat\A}(\Omega^\star)-S_{\widehat\A}(\Omega)
    \ge
    S_\A(\Omega^\star)-S_\A(\Omega)-2\delta_\A
    \ge
    \Delta_\A-2\delta_\A.
\]
If empirical deviation is less than $(\Delta_\A-2\delta_\A)/2$ uniformly, no
non-equivalent candidate can beat the true class. Apply the concentration
assumption with this value.
\end{proof}

\section{Intervention Selection as Handle-Operation Planning}

Given posterior
\[
    P(\Omega,K,\Gamma,Q\mid \D),
\]
a test is not an abstract $\doop(X_i=x)$ but a handle-operation pair $(h,u)$.
Its value is expected posterior contraction.

\begin{definition}[Expected information gain of a handle-operation]
For $u\in\U_h$,
\begin{equation}
    \EIG(h,u)
    =
    \Expect_{y'\sim P(\cdot\mid \D,h,u)}
    \left[
        \KL\!\left(
            P(\Omega\mid \D,h,u,y')
            \,\|\, 
            P(\Omega\mid \D)
        \right)
    \right].
    \label{eq:eig}
\end{equation}
\end{definition}

The recommended test is
\begin{equation}
    (h^\star,u^\star)
    =
    \argmax_{h,u}
    \left[
        \EIG(h,u)
        -
        \lambda_c\cost(h,u)
        -
        \lambda_r\risk(h,u)
        -
        \lambda_\chi \mathbf 1\{(h,u)\notin\mathcal C\}
    \right].
    \label{eq:best-intervention}
\end{equation}

\begin{proposition}[Diagnostic separation form]
Suppose two hypotheses $\Omega_1,\Omega_2$ dominate the posterior. Let
\[
    p_i^{h,u}(y')
    =
    P(y'\mid \Omega_i,\D,h,u).
\]
Then, up to posterior weights, $\EIG(h,u)$ is increasing in the predictive
separation
\[
    \KL(p_1^{h,u}\|p_2^{h,u})
    +
    \KL(p_2^{h,u}\|p_1^{h,u}).
\]
Thus good interventions are handle-operations for which competing loop
hypotheses make different downstream predictions.
\end{proposition}

\section{Agency Tests as Handle Tests}

The usual intuitive agency tests become precise only when indexed by handles.
Table~\ref{tab:handle-tests} gives the translation.

\begin{table}[h]
\centering
\small
\begin{tabular}{@{}p{0.17\linewidth}p{0.33\linewidth}p{0.42\linewidth}@{}}
\toprule
Suspected structure & Handle operation & Diagnostic contrast \\
\midrule
Sensor channel $Q^S$ &
hide, delay, corrupt, amplify input handle &
behavior tracks information rather than world state \\
Action channel $Q^A$ &
block, attenuate, reroute output handle &
environmental control collapses or compensates \\
Goal/resource $G$ &
move, remove, revalue resource handle &
means change while latent target persists \\
Belief state $B$ &
false cue, occlusion, asymmetric information &
action follows belief rather than actual state \\
Policy $\pi$ &
change cost, risk, payoff, time pressure &
action trades off costs approximately coherently \\
Boundary $C$ &
separate, perturb, damage, constrain coupling &
repair, compensation, or loss of autonomy \\
Communication &
cut, corrupt, delay signal handle &
coordination degrades predictably \\
Model of other agent &
alter partner behavior or apparent belief &
anticipation, teaching, deception, or adaptation \\
Self-maintenance &
resource stress, load, disruption &
actions preserve viability variables \\
\bottomrule
\end{tabular}
\caption{Agency tests rephrased as embedded handle operations.}
\label{tab:handle-tests}
\end{table}

\section{Language as a High-Bandwidth Handle}

Language is special because it makes latent control variables queryable. A
non-linguistic observer sees trajectories:
\[
    Q^S_t,Q^A_t,Q^S_{t+1},Q^A_{t+1},\ldots
\]
and must infer $B_t,G_t,\pi$. A linguistic system emits public symbols that
purport to describe these latents:
\[
    B_t \mapsto \ulcorner B_t\urcorner,\qquad
    G_t \mapsto \ulcorner G_t\urcorner,\qquad
    \pi_t \mapsto \ulcorner \pi_t\urcorner .
\]
This does not guarantee truth. In adversarial settings, self-reports may be
strategic. But language changes the access model by adding handles:
\[
    h_{\mathrm{ask}},\quad
    h_{\mathrm{challenge}},\quad
    h_{\mathrm{explain}},\quad
    h_{\mathrm{commit}},\quad
    h_{\mathrm{counterfactual}}.
\]
These handles can generate high expected information gain when the system is
cooperative or only weakly adversarial.

\section{Algorithmic Feeds as Hidden Agent Loops}

An algorithmic feed illustrates handle failure. A user observes:
\[
    \text{content item}\to \text{reaction}\to \text{next content item}.
\]
But the relevant loop is closer to:
\[
    \text{platform objective}
    \to
    \text{ranking policy}
    \to
    \text{user-state perturbation}
    \to
    \text{telemetry}
    \to
    \text{policy update}.
\]
The user lacks handles for:
\[
    G_{\mathrm{platform}},
    \quad
    Q^A_{\mathrm{ranking}},
    \quad
    Q^S_{\mathrm{telemetry}},
    \quad
    \pi_{\mathrm{update}},
    \quad
    \text{auction and advertiser constraints}.
\]
Thus user-level observations can be explained as local content variation even
when a distributed optimization loop is acting on the user. Agency is hidden not
because the loop is weak, but because the access model is poor.

A meaningful audit interface would expose at least:
\begin{align*}
    &\text{candidate set},\quad
    \text{ranking action},\quad
    \text{objective contribution}, \\
    &\text{targeting features},\quad
    \text{measured response},\quad
    \text{policy update}.
\end{align*}
Such an interface changes $K$ and $\Gamma$, increasing the recoverable margin
$\Delta_\A$ and shrinking access-equivalence classes.

\section{Data Contracts}

Access-model UAD requires different claims for different data tiers.

\begin{definition}[Tiered data contract]
\begin{enumerate}[label=\textbf{Tier \arabic*:},leftmargin=*]
    \item Passive access:
    \[
        \D=(Y_{1:T},K).
    \]
    The system can generate candidate loops but cannot make strong causal claims.

    \item Manipulable access:
    \[
        \D=(Y_{1:T},K,\Hh^U,\Gamma,\cost,\risk).
    \]
    The system can recommend handle-operations and update loop posteriors from
    their outcomes.

    \item Interface-exposed access:
    \[
        \D=(Q^S_t,Q^A_t,R_t,\text{update}_t,Y_{t+1})_{t=1}^T.
    \]
    The system can audit a largely exposed control loop.
\end{enumerate}
\end{definition}

The same mathematical model can operate at all tiers, but the strength of
conclusions differs:
\[
    \text{passive discovery}
    <
    \text{causal testing}
    <
    \text{loop auditing}.
\]

\section{Practical Algorithm}

\begin{enumerate}[leftmargin=*]
    \item \textbf{Fit access model.}
    Estimate or encode observation kernels $K_h$, operation kernels $\Gamma_h$,
    costs, risks, legal constraints, and handle semantics.

    \item \textbf{Learn representations.}
    Produce multi-scale representations $Z^{\sigma,\tau}$ through coarse-graining,
    prediction, and event segmentation.

    \item \textbf{Propose candidates.}
    Use amortized heads to propose
    \[
        q_\phi(C,Q,B,G,\pi,M\mid \D).
    \]

    \item \textbf{Score candidates.}
    Evaluate explicit UAD/EIS/B-IQ-style energies~\cite{uad2025,eis2025,biq2025}:
    \[
        J_Q(C),\quad
        I_{\mathrm{ctrl}},\quad
        R_{\mathrm{persist}},\quad
        R_{\mathrm{belief}},\quad
        \Delta L_{\mathrm{EIS}},\quad
        \DL.
    \]

    \item \textbf{Infer posterior.}
    Approximate
    \[
        P(\Omega,Q,K,\Gamma\mid \D)
    \]
    by variational inference, MCMC, sequential Monte Carlo, or energy-based
    reranking.

    \item \textbf{Select test.}
    Choose $(h,u)$ maximizing Eq.~\eqref{eq:best-intervention}.

    \item \textbf{Update.}
    Observe outcome $y'$ and update posterior. Stop when posterior mass over
    agent-loop equivalence classes is sufficiently concentrated or when further
    tests are too costly.
\end{enumerate}

\section{Discussion}

The handle approach changes the ontology of intervention-capable UAD. It says:
do not start with variables and interventions. Start with access. Variables,
interfaces, interventions, and intentions are all high-level abstractions over
embedded dynamics.

This has several consequences.

First, agency claims become conditional:
\[
    \text{``agentic under access model } \A\text{''}.
\]
No method can recover distinctions erased by $\A$.

Second, intervention suggestion becomes scientifically meaningful only when
operation semantics are modeled. A recommendation such as ``perturb variable
$X_i$'' is often ill-posed. A recommendation such as ``use handle $h$ with
operation $u$ because hypotheses $\Omega_1$ and $\Omega_2$ predict different
outcomes'' is well-formed.

Third, interface forcing is access-channel design. It may improve identifiability,
but it may also distort the system. The correct question is whether the forced
interface is faithful enough for the target agency claim.

Fourth, layer-wise regularization should be weakened. Lower layers should be
regularized toward stable, local, separable, and manipulable roles. Full agent-loop
regularization should be sparse and scale-selective.

Fifth, adversarial systems break naive EIS~\cite{eis2025,dennett1987intentional}. If a system optimizes its outputs to
imitate the sufficient statistics of agency while hiding its loop, behavioral
compression alone is insufficient. One needs handle diversity, hidden tests,
costly perturbations, or enforced audit interfaces.

\section{Conclusion}

Intervention, interface, and intention are all compressed names for deep embedded
stacks. Causal graphs make intervention look primitive; real systems provide
handles. Intentional explanations make goals look primitive; real agents expose
or hide them through behavior, language, and interface structure. UAD~\cite{uad2025} should
therefore be formulated relative to an access model:
\[
    \A=(\Hh^Y,\Hh^U,K,\Gamma,\cost,\risk,\mathcal C).
\]
Given $\A$, the target of discovery is not an agent simpliciter but a posterior
over handle-mediated loop hypotheses:
\[
    P(C,Q,B,G,\pi,K,\Gamma\mid \D).
\]
The recoverable object is an access-equivalence class. The appropriate test is
not $\doop(X_i=x)$ but a handle-operation chosen by expected information gain.
The practical path is therefore:
\begin{align*}
    &\text{access modeling}
    \to \text{role-regularized representation learning} \\
    &\to \text{sparse agent-loop inference}
    \to \text{handle-based experimental design}.
\end{align*}
This makes UAD less elegant but more real. It also explains why hidden optimization
systems such as algorithmic feeds are hard for humans to perceive: they act through
the handles we see while hiding the handles that would reveal the agent loop.

\bibliographystyle{IEEEtran}
\bibliography{refs}

\end{document}
